\section{Fazit und Ausblick}

\subsection{Zusammenfassung}

Ausgangspunkt der vorliegenden Arbeit war die Beobachtung, dass Reinforcement-Learning (RL)-Agenten im Schienenverkehr ohne eine geeignete Initialisierung kaum konvergieren und bei einer rein simulativen Offline-Initialisierung zudem das Risiko des Overfittings auf spezifische Trainingsstrecken besteht. Als Reaktion darauf verfolgte diese Arbeit das Ziel, mittels Behavior Cloning (BC) ein neuronales Netz zu entwickeln, das reale Fahrstrategien menschlicher Triebfahrzeugführer aus historischen Betriebsdaten erlernt und damit als robuste, generalisierbare Baseline für einen zukünftigen Warm-Start von RL-Agenten dienen kann.

Zur Erreichung dieser Zielsetzung wurde ein iteratives, experimentelles Vorgehen gewählt. Nach einer theoretischen Fundierung zu Supervised Learning, BC und geeigneten Netzwerkarchitekturen folgte eine ausführliche Exploratory Data Analysis der realen Betriebsdaten, auf deren Grundlage eine deterministische Preprocessing- und Feature-Engineering-Pipeline entwickelt wurde. Über ein fahrtenbasiertes Trip-Level-Splitting wurde sichergestellt, dass Trainings-, Validierungs- und Testmenge vollständig disjunkte, unabhängige Fahrten enthalten und Data Leakage auf Zeitschrittebene ausgeschlossen ist. In einem dreistufigen Auswahlprozess wurden anschließend Architektur, Netzwerkgröße sowie Optimierer und Verlustfunktion empirisch bestimmt; als finale Konfiguration ergab sich ein LSTM mit der Layer-Konfiguration 128-64-256-128 (650.369 Parameter), trainiert mit Adam und Smooth-L1-Loss. Über fünf Trainingsiterationen wurden Rauschen und Gewichtsinitialisierung, Batch-Size sowie zyklische Regularisierungstechniken systematisch angepasst, bevor das finale Modell auf dem vollständigen Trainingsdatensatz trainiert wurde.

Die abschließende Evaluation auf der bis dahin unangetasteten Testmenge zeichnet ein gemischtes Bild: Auf den unmittelbar über den quadratischen Fehler definierten Metriken ($R^2$, RMSE, MSE) übertrifft das Modell eine strenge Persistenz-Baseline deutlich, während es bei MAE sowie bei Richtungs- und Toleranz-Accuracy hinter dieser zurückbleibt. Die klassenweise Analyse identifiziert die Ursache: Seltene, aber sicherheitsrelevante Bremsanforderungen werden mit einem Recall von lediglich 3,2\,\% nahezu systematisch verpasst, während das Modell in den weitaus häufigeren, ruhigen Beharrungsphasen präzise Vorhersagen liefert. Dieses Verhalten lässt sich plausibel auf die begrenzte Datenbasis von 1.388 Trips bzw. rund 178 Stunden Aufzeichnungszeit zurückführen, in der Extremmanöver naturgemäß unterrepräsentiert sind. Insgesamt liefert das entwickelte Modell damit keine im vollen Umfang optimalen, aber für den angestrebten Machbarkeitsnachweis hinreichende Ergebnisse.

\subsection{Beantwortung der Forschungsfrage}
\label{subsec:forschungsfrage}

Um die eingangs formulierte Forschungsfrage fundiert beantworten zu können, wird zunächst auf die Erreichung der drei in Abschnitt~\ref{subsec:zielsetzung} definierten Teilziele eingegangen, bevor eine zusammenfassende Bewertung der Forschungsfrage selbst erfolgt.

Das erste Teilziel, die Transformation der unstrukturierten Rohdaten in eine strukturierte, für das maschinelle Lernen nutzbare Repräsentation, wurde vollständig erreicht. Mittels der in Kapitel~\ref{sec:daten_preprocessing} entwickelten EDA- und Feature-Engineering-Pipeline konnten zentrale, fahrphysikalisch kausale Input-Features wie Streckengradient, Geschwindigkeit und Distanz zu Stoppstellen identifiziert und in einen wohldefinierten Zustands- sowie Aktionsraum überführt werden. Die dabei entstandene deterministische Pipeline stellt zugleich sicher, dass Preprocessing und Feature-Auswahl für alle nachfolgenden Experimentierschritte reproduzierbar bleiben.

Auch das zweite Teilziel, die Konzeptionierung und Implementierung einer geeigneten Deep-Learning-Architektur, wurde erreicht. Der in Kapitel~\ref{sec:konzeption} durchgeführte dreistufige Auswahlprozess führte zu einer LSTM-Architektur, die in der Lage ist, die nicht-linearen, zeitlich abhängigen Zusammenhänge zwischen den topologischen Streckendaten und der kontinuierlichen Zielgröße abzubilden. Sowohl die Netzwerkdimensionierung als auch die Wahl von Optimierer und Verlustfunktion wurden dabei empirisch abgesichert und nicht willkürlich festgelegt.

Das dritte Teilziel, die Verhinderung von Overfitting bei gleichzeitiger Sicherstellung der Generalisierungsfähigkeit, wurde nur teilweise erreicht. Durch das fahrtenbasierte Trip-Level-Splitting sowie die in Kapitel~\ref{sec:training} angewandten Regularisierungstechniken generalisiert das Modell nachweislich auf vollständig ungesehene Fahrten und übertrifft die Persistenz-Baseline bei den regressionsbasierten Metriken. Gleichzeitig zeigt die in Kapitel~\ref{sec:evaluation} durchgeführte klassenweise Analyse, dass diese Generalisierung nicht gleichmäßig über alle Fahrsituationen hinweg gelingt: Insbesondere für seltene, dynamische Brems- und Beschleunigungsmanöver bleibt die Generalisierungsfähigkeit deutlich hinter der für ruhige Beharrungsphasen erreichten Güte zurück.

Auf dieser Grundlage lässt sich die zentrale Forschungsfrage, ob sich aus unstrukturierten historischen Betriebsdaten ein künstliches neuronales Netz trainieren lässt, das eine generalisierbare Fahrstrategie für Züge erlernt, mit einem begründeten Ja beantworten. Der im Rahmen dieser Arbeit erbrachte Machbarkeitsnachweis zeigt, dass ein solches Training grundsätzlich möglich ist: Das Modell erlernt aus realen, unstrukturierten Sensordaten eine Policy, die auf ungesehenen Testfahrten insgesamt bessere Regressionsergebnisse liefert als eine triviale Fortschreibungsstrategie und in ruhigen Fahrsituationen sehr präzise Vorhersagen trifft. Zugleich macht die Arbeit jedoch deutlich, dass die Qualität einer solchen Policy unmittelbar von der zur Verfügung stehenden Datenmenge abhängt: Gerade die für eine sichere und vollständig generalisierbare Fahrstrategie entscheidenden, aber seltenen Extremmanöver konnten aufgrund der vergleichsweise kleinen Datenbasis von 178 Stunden nicht robust erlernt werden. Dies ist zugleich der zentrale Grund dafür, dass die erzielten Ergebnisse am Ende dieser Arbeit nicht als optimal, wohl aber als für den angestrebten Machbarkeitsnachweis ausreichend zu bewerten sind. Die Forschungsfrage ist damit im Grundsatz positiv, aber mit der wesentlichen Einschränkung einer ausreichend großen und hinsichtlich seltener Ereignisse hinreichend reichhaltigen Datenbasis zu beantworten.

\subsection{Kritische Reflexion}

Trotz des grundsätzlich positiven Machbarkeitsnachweises sind mehrere Aspekte des gewählten Vorgehens kritisch zu hinterfragen.

Zunächst betrifft dies die Zusammensetzung der Datenbasis selbst. Da sämtliche Trips von einer einzigen, homogenen Güterverkehrsstrecke stammen, konnte kein streckenübergreifendes Splitting vorgenommen werden; die in Abschnitt~\ref{subsec:finale_testevaluation} beschriebene Dominanz von Rangier- und Anfahrsituationen im Testset ist eine unmittelbare Folge dieser Datengrundlage. Die berichteten Metriken sind somit zunächst nur für diese spezifische Strecke und das zugrundeliegende Fahrzeug- bzw. Führerprofil aussagekräftig; inwieweit die erlernte Policy auf andere Streckentypen, Zuggattungen oder Traktionsarten übertragbar ist, kann im Rahmen dieser Arbeit nicht beurteilt werden.

Methodisch ist zudem die ausgeprägte Klassenimbalance zwischen Halte-, Brems- und Beschleunigungsanforderungen zu nennen. Die verwendete Smooth-L1-Verlustfunktion optimiert den Gesamtfehler primär über die zahlenmäßig dominierenden, unauffälligen Zeitschritte, wodurch seltene, aber sicherheitsrelevante Ereignisse systematisch unterrepräsentiert bleiben. Gezielte Gegenmaßnahmen wie eine klassengewichtete Verlustfunktion oder ein gezieltes Oversampling dynamischer Fahrabschnitte wurden in dieser Arbeit bewusst nicht eingesetzt, um die Vergleichbarkeit der in Kapitel~\ref{sec:konzeption} und~\ref{sec:training} durchgeführten Experimente nicht zusätzlich zu verkomplizieren; ihr Fehlen dürfte jedoch zu der beobachteten, unzureichenden Erkennung von Bremsanforderungen beigetragen haben.

Auch der gewählte methodische Rahmen des Behavior Cloning selbst ist limitiert: Da das Modell rein offline auf historischen Zustands-Aktions-Paaren trainiert wurde, ohne jemals mit den Konsequenzen eigener Entscheidungen in einer geschlossenen Regelschleife konfrontiert zu sein, bleibt unklar, wie sich kleine, im Training unvermeidbare Abweichungen von der menschlichen Referenz im tatsächlichen Einsatz akkumulieren würden. Die in Abschnitt~\ref{sec:problemstellung} beschriebene Gefahr des Extrapolationsfehlers bei ungesehenen Situationen wurde durch das Trip-Level-Splitting zwar für die Evaluation, nicht aber für einen realen, sequenziellen Einsatz methodisch adressiert.

Schließlich wurden die Trainings-Hyperparameter Lernrate, Batch-Size, Weight Decay und Dropout-Rate während der Architektur- und Dimensionierungsexperimente in Kapitel~\ref{sec:konzeption} bewusst auf feste Startwerte gesetzt, um die Vergleichbarkeit der jeweils untersuchten Designentscheidung sicherzustellen. Dies erhöht zwar die interne Validität der einzelnen Experimentstufen, schließt jedoch nicht aus, dass eine gemeinsame, gekoppelte Optimierung von Architektur und Hyperparametern zu einer anderen finalen Konfiguration und potenziell besseren Ergebnissen geführt hätte.

\subsection{Ausblick}

Ausgehend von der in Abschnitt~\ref{subsec:kritische_einordnung} sowie Abschnitt~\ref{subsec:forschungsfrage} identifizierten Kernursache der beobachteten Schwächen ergibt sich als naheliegendste Anschlussarbeit die Erweiterung der zugrunde liegenden Datenbasis. Da sich die Modellgüte gerade bei seltenen Brems- und Beschleunigungsmanövern unmittelbar an der geringen Anzahl entsprechender Trainingsbeispiele festmachen lässt, würde bereits eine Verdopplung oder Verdreifachung der aufgezeichneten Betriebsstunden die absolute Zahl solcher Extremereignisse proportional erhöhen und damit voraussichtlich zu einer spürbar robusteren Erkennung dynamischer Fahrsituationen führen. Ergänzend dazu erscheinen datengetriebene Maßnahmen wie eine klassengewichtete Verlustfunktion oder ein gezieltes Oversampling unterrepräsentierter Manöver als vielversprechende, vergleichsweise aufwandsarme nächste Schritte.

Der zentrale nächste Schritt liegt jedoch, wie bereits in der Zielsetzung dieser Arbeit angelegt, in der eigentlichen Nutzung des entwickelten Modells als Warm-Start-Initialisierung für einen Reinforcement-Learning-Agenten. Anstatt einen RL-Agenten wie in Abschnitt~\ref{sec:problemstellung} beschrieben ohne fachliche Vorprägung und mit dem damit verbundenen Risiko einer mangelnden Konvergenz zu initialisieren, kann die in dieser Arbeit trainierte BC-Policy als Startpunkt dienen, von dem aus im Anschluss ein RL-Training durchgeführt wird. Auf diesem Weg ließe sich untersuchen, ob und in welchem Umfang sich das Fahrverhalten des Agenten durch anschließendes Reinforcement Learning gegenüber der reinen BC-Baseline tatsächlich verbessern lässt – insbesondere in genau jenen dynamischen Übergangsphasen, in denen das in dieser Arbeit entwickelte Modell noch die größten Schwächen zeigt. Erst eine solche kombinierte Warm-Start-plus-RL-Strategie würde damit das übergeordnete, in Abschnitt~\ref{subsec:zielsetzung} formulierte Ziel dieser Arbeit vollständig einlösen: eine zuverlässige, generalisierbare Fahrstrategie für den automatisierten Schienenverkehr bereitzustellen.